\chapterpage{Sorting, intervals and comparators}{Seven problems. Writing a
correct multi-key comparator under time pressure is a skill in itself, and
interval sweeps are the most common ``medium'' in a screen.}

% ------------------------------------------------------------------
\problem{Sort by frequency, then by value}{Medium}

Sort the array so that more frequent values come first, and equally frequent
values come in ascending numeric order.
\texttt{[1,1,2,2,2,3]} gives \texttt{[2,2,2,1,1,3]}.

\idea{Count first, then sort the \emph{boxed} values with a two-key comparator.
Primitive arrays cannot take a comparator, so boxing is required.}

\solheading
\begin{lstlisting}[style=java]
public static int[] sortByFrequency(int[] a) {
    Map<Integer,Integer> freq = new HashMap<>();
    for (int x : a) freq.merge(x, 1, Integer::sum);

    Integer[] boxed = new Integer[a.length];
    for (int i = 0; i < a.length; i++) boxed[i] = a[i];

    Arrays.sort(boxed, (x, y) -> {
        int fx = freq.get(x), fy = freq.get(y);
        if (fx != fy) return Integer.compare(fy, fx);   // frequency desc
        return Integer.compare(x, y);                   // value asc
    });

    int[] out = new int[a.length];
    for (int i = 0; i < a.length; i++) out[i] = boxed[i];
    return out;
}
\end{lstlisting}

\cost{$O(n\log n)$ time, $O(n)$ space.}
\trap{\texttt{Arrays.sort(int[], Comparator)} does not exist --- there is no
such overload, and the compile error surprises people under pressure. Box to
\texttt{Integer[]} or sort a \texttt{List<Integer>}.}

% ------------------------------------------------------------------
\problem{Merge overlapping intervals}{Medium}

Given intervals as \texttt{[start, end]}, merge every pair that overlaps.
\texttt{[[1,3],[2,6],[8,10]]} gives \texttt{[[1,6],[8,10]]}.

\idea{Sort by start. Then one pass: either the next interval starts inside the
one you are building, in which case stretch its end, or it does not, in which
case close the current one and open a new one.}

\solheading
\begin{lstlisting}[style=java]
public static int[][] merge(int[][] intervals) {
    if (intervals == null || intervals.length == 0) return new int[0][];
    int[][] a = intervals.clone();
    Arrays.sort(a, (x, y) -> Integer.compare(x[0], y[0]));

    List<int[]> out = new ArrayList<>();
    int[] current = new int[] {a[0][0], a[0][1]};
    for (int i = 1; i < a.length; i++) {
        if (a[i][0] <= current[1]) {
            current[1] = Math.max(current[1], a[i][1]);
        } else {
            out.add(current);
            current = new int[] {a[i][0], a[i][1]};
        }
    }
    out.add(current);
    return out.toArray(new int[0][]);
}
\end{lstlisting}

\cost{$O(n\log n)$ time, $O(n)$ space.}
\trap{\texttt{current[1] = a[i][1]} without the \texttt{Math.max} loses a
fully contained interval: merging \texttt{[1,10]} with \texttt{[2,3]} must
still end at 10. Also remember to add the last \texttt{current} after the
loop.}

% ------------------------------------------------------------------
\problem{Can one person attend all meetings?}{Easy}

Given meeting intervals, return \texttt{true} if none of them overlap.

\idea{Sort by start; any overlap is then between neighbours.}

\solheading
\begin{lstlisting}[style=java]
public static boolean canAttendAll(int[][] meetings) {
    if (meetings == null || meetings.length < 2) return true;
    int[][] a = meetings.clone();
    Arrays.sort(a, (x, y) -> Integer.compare(x[0], y[0]));
    for (int i = 1; i < a.length; i++)
        if (a[i][0] < a[i - 1][1]) return false;
    return true;
}
\end{lstlisting}

\cost{$O(n\log n)$ time.}
\trap{Whether touching intervals count as overlapping. \texttt{[1,5]} and
\texttt{[5,9]} do not overlap if the end is exclusive --- hence \texttt{<} and
not \texttt{<=}. The statement decides; read it.}

% ------------------------------------------------------------------
\problem{Minimum number of meeting rooms}{Medium}

Return the largest number of meetings running at the same time.

\idea{A sweep. Put every start and every end on a timeline, sort both, then walk
them: a start raises the count, an end lowers it. Track the peak. Ends must
be processed before starts at the same instant, since a room freed at 5 can
be reused at 5.}

\solheading
\begin{lstlisting}[style=java]
public static int minRooms(int[][] meetings) {
    if (meetings == null || meetings.length == 0) return 0;
    int n = meetings.length;
    int[] starts = new int[n], ends = new int[n];
    for (int i = 0; i < n; i++) {
        starts[i] = meetings[i][0];
        ends[i]   = meetings[i][1];
    }
    Arrays.sort(starts);
    Arrays.sort(ends);

    int rooms = 0, peak = 0, e = 0;
    for (int s = 0; s < n; s++) {
        while (e < n && ends[e] <= starts[s]) { rooms--; e++; }
        rooms++;
        peak = Math.max(peak, rooms);
    }
    return peak;
}
\end{lstlisting}

\cost{$O(n\log n)$ time, $O(n)$ space.}
\trap{\texttt{ends[e] < starts[s]} instead of \texttt{<=} books a second room
for back-to-back meetings. The two arrays are sorted independently on
purpose --- the pairing between a particular start and end does not matter.}

% ------------------------------------------------------------------
\problem{Insert an interval}{Medium}

The intervals are already sorted and non-overlapping. Insert a new one, merging
where it overlaps.

\idea{Three phases, no sorting needed: copy everything that ends before the new
one starts, absorb everything that overlaps into a widened interval, then
copy the rest.}

\solheading
\begin{lstlisting}[style=java]
public static int[][] insert(int[][] intervals, int[] newInterval) {
    List<int[]> out = new ArrayList<>();
    int i = 0, n = intervals.length;
    int start = newInterval[0], end = newInterval[1];

    while (i < n && intervals[i][1] < start) out.add(intervals[i++]);

    while (i < n && intervals[i][0] <= end) {
        start = Math.min(start, intervals[i][0]);
        end   = Math.max(end,   intervals[i][1]);
        i++;
    }
    out.add(new int[] {start, end});

    while (i < n) out.add(intervals[i++]);
    return out.toArray(new int[0][]);
}
\end{lstlisting}

\cost{$O(n)$ time --- the input is already sorted, so do not re-sort it.}
\trap{Adding the merged interval inside the second loop, which emits it once
per overlap. It goes in exactly once, after the loop.}

% ------------------------------------------------------------------
\problem{Sort employees by salary then name}{Easy}

Sort records by salary descending, and alphabetically for equal salaries. A
pure comparator drill --- this exact shape appears in data-processing tasks.

\idea{The \texttt{Comparator} factory methods compose. Build the ordering
declaratively rather than hand-rolling the branches.}

\solheading
\begin{lstlisting}[style=java]
record Employee(String name, int salary) {}

public static void sortEmployees(List<Employee> staff) {
    staff.sort(Comparator.comparingInt(Employee::salary).reversed()
                         .thenComparing(Employee::name));
}

// the hand-written equivalent, when you cannot remember the fluent form
public static void sortEmployeesManually(List<Employee> staff) {
    staff.sort((a, b) -> {
        if (a.salary() != b.salary())
            return Integer.compare(b.salary(), a.salary());
        return a.name().compareTo(b.name());
    });
}
\end{lstlisting}

\cost{$O(n\log n)$.}
\trap{\texttt{.reversed()} reverses \emph{everything composed so far}. Chaining
it on after \texttt{thenComparing(name)} reverses the name order as well. Put
\texttt{reversed()} on the salary key alone, as above.}

% ------------------------------------------------------------------
\problem{Kth largest element}{Medium}

Return the $k$th largest value in an unsorted array (the $k$th in sorted
descending order, duplicates counted).

\idea{Keep a min-heap of size $k$. The smallest thing in it is, once you have
seen everything, the $k$th largest overall.}

\solheading
\begin{lstlisting}[style=java]
public static int kthLargest(int[] a, int k) {
    PriorityQueue<Integer> heap = new PriorityQueue<>();   // min-heap
    for (int x : a) {
        heap.offer(x);
        if (heap.size() > k) heap.poll();
    }
    return heap.peek();
}
\end{lstlisting}

\cost{$O(n\log k)$ time, $O(k)$ space. Sorting is $O(n\log n)$ and perfectly
acceptable when $k$ is close to $n$.}
\trap{Using a max-heap and polling $k$ times is $O(n + k\log n)$ to build and
drain --- fine, but it needs all $n$ elements in memory. The size-$k$ min-heap
is the answer when the input is a stream.}
